%
% patch_atomicity.tex
%
% A Z specification of the lux Element-ABC apply_patch boundary: an atomic,
% all-or-nothing multi-field patch on a single element whose selection AND rows
% writes are observed by a composed Hub-side model (the filtered table's full
% selection and dataset).  Formalises the rollback-atomicity discipline hardened
% across three fix/review rounds -- a composed mutable model that the shallow
% vars() snapshot could not restore; an observer notified INSIDE the patch that a
% later failing key could not un-fire; and, added here, a REENTRANT commit-phase
% notification whose observer folds a rows refresh into the model and re-projects
% by issuing a NESTED apply_patch on the same element.  ProB enumerates every
% patch shape and nesting order and either confirms the atomicity invariant or
% produces the exact offending trace.
%
% Source of record:
%   src/punt_lux/domain/element_abc.py:202  -- Element.apply_patch: snapshot =
%                                              dict(vars(self)); open a per-patch
%                                              _notify_buffer (reset to None on
%                                              entry, so each nested apply_patch
%                                              owns its own buffer); run _set_<key>
%                                              per entry; on any raise restore
%                                              vars(self), discard the buffer, and
%                                              re-raise; on success set the buffer
%                                              to None and THEN fire each deferred
%                                              notification -- so notifications
%                                              fired during commit land in the
%                                              nested patch's fresh buffer, not the
%                                              outer one, which is already None.
%   src/punt_lux/protocol/elements/table.py:180 -- _set_rows: reconciles the
%                                              selection and notifies "rows" (an
%                                              external dataset refresh) plus,
%                                              on a reconcile, "selected_row_ids".
%   src/punt_lux/protocol/elements/table.py:213 -- _set_selected_row_ids:
%                                              notifies "selected_row_ids".
%   src/punt_lux/protocol/compositions/filtered_table_model.py:118 --
%                                              _on_table_change: the observer.  On
%                                              "selected_row_ids" it folds the
%                                              write into the full selection.  On
%                                              "rows", if NOT self._reprojecting,
%                                              it ABSORBS the write as the new
%                                              dataset and _reproject()s, which
%                                              sets self._reprojecting = True and
%                                              issues a NESTED apply_patch writing
%                                              the filtered subset + projected
%                                              selection; the guard excludes that
%                                              inner write from being re-absorbed.
%
% Type-check:  fuzz -t docs/patch_atomicity.tex
% Model-check (see Section "Verification" for the exact goal predicates):
%   probcli docs/patch_atomicity.tex -model_check                 (deadlock-freedom)
%   probcli docs/patch_atomicity.tex -model_check -nodead \
%     -goal "size(smode)=0 & (observed(selObs)/=fields(ksel) or
%            observed(theModel)/=extern)"                          (invariant AT)
%
% Formatting follows the fuzz house rules: \quad~ for continuation indent
% (never \t1..\t9), schema lines under 80 columns, predicates broken at
% \land/\lor with the operator starting the continuation line.
%

\documentclass[a4paper,10pt,fleqn]{article}
\usepackage[margin=1in]{geometry}
\usepackage{fuzz}
\usepackage[colorlinks=true,linkcolor=blue,citecolor=blue,urlcolor=blue]{hyperref}

\begin{document}

\title{The lux \texttt{apply\_patch} Atomicity Boundary under Reentrant
Notification: a Z Specification}
\author{Formal model of the all-or-nothing patch discipline with nested patches}
\date{July 2026}
\maketitle

% ============================================================
\section{Introduction}
% ============================================================

An agent patches a single Element-ABC element by sending a mapping of field
names to new values.  \texttt{Element.apply\_patch} applies the entries in
order, each through a \texttt{\_set\_<key>} setter, and promises
\emph{all-or-nothing}: if any setter rejects its value, every field the patch
already touched is restored and the element is left exactly as it was before
the patch began.  A composed table also has an \emph{observer} --- the filtered
table's Hub-side model --- that folds a selection write into a full selection
spanning filter-hidden rows, and folds an external \emph{rows} write in as its
new dataset.  The property that must hold is that \emph{a rejected patch
produces no observable change at all}, and that a patch that succeeds leaves
every observer reflecting the patch's final values.

The same class of defect --- a rejected patch leaving an observable trace, or a
committed patch leaving the model inconsistent --- recurred across three
fix/review rounds:

\begin{enumerate}
  \item the table's selection lived in a \emph{mutable} composed model, so the
    shallow \texttt{dict(vars(self))} snapshot captured the model
    \emph{reference}; a setter mutated that same object in place, and the
    rollback restored a reference to an already-mutated model;
  \item the selection setter notified observers \emph{inside} the patch, so a
    composed model folded the intermediate selection before a later key was
    rejected; the rollback restored the element's fields but could not un-fire
    the notification;
  \item the model absorbed an external rows refresh as its dataset and
    re-projected by issuing a \emph{nested} \texttt{apply\_patch} from inside
    the outer patch's commit-phase notification.  Without a guard the model
    folds \emph{its own} filtered subset back in as the dataset, and the nested
    write re-triggers the observer without bound.
\end{enumerate}

Round~1 was closed by making the composed selection model immutable.  Round~2
by \emph{deferring} observer notification to commit.  Round~3 is closed by two
mechanisms working together: each \texttt{apply\_patch} call owns its own
notification buffer (it is reset to \texttt{None} on entry, so a notification
fired during an outer commit lands in the nested patch's fresh buffer, never the
outer one), and a \texttt{\_reprojecting} boolean guard on the model excludes
the model's own inner write from re-absorption.  All three fixes survive here as
fidelity variants (Section~\ref{sec:fidelity}) that reintroduce each defect to
confirm the model still detects it.  Empirical round-by-round fixing kept
missing an interleaving, so this document models the boundary as a finite state
machine over a bounded \emph{stack} of patch frames, letting ProB enumerate
every patch shape and every nesting order exhaustively.

% ============================================================
\section{Basic Types}
% ============================================================

Because the reentrancy is tied to specific fields (a \texttt{rows} write
triggers a re-projection; a \texttt{selected\_row\_ids} write only folds) and to
one distinguished observer (the model), the carriers are modelled concretely as
free types rather than as unconstrained given sets.  This bakes in exactly the
topology under test --- the model watches \texttt{rows} and re-projects, a
plain observer watches the selection --- and keeps every carrier finite, so ProB
is exhaustive without a \texttt{DEFAULT\_SETSIZE} parameter.

\begin{zed}
  FMODE ::= setting | flushing
\end{zed}

A patch frame is \texttt{setting} while its setters run (buffering
notifications) and \texttt{flushing} while its committed notifications are
delivered.

\begin{zed}
  SWITCH ::= off | on
\end{zed}

Used for the \texttt{\_reprojecting} guard and for a frame's ``is a
re-projection'' marker.

\begin{zed}
  GOODVAL ::= a | b
\end{zed}

The values a field may validly hold.  Two distinct values are enough to show a
field \emph{changing} and to distinguish an agent's pushed dataset (modelled as
\texttt{a}) from the filtered subset a re-projection writes back (modelled as
\texttt{b}), so that ``the model absorbed its own subset'' is observable.

\begin{zed}
  KEY ::= krows | ksel
\end{zed}

The two fields under test: \texttt{krows} (the dataset) and \texttt{ksel} (the
selection).

\begin{zed}
  OBS ::= theModel | selObs
\end{zed}

Two observers: \texttt{theModel} is the filtered table's model, which watches
\texttt{krows} and re-projects; \texttt{selObs} is a plain observer of the
selection, which only folds.

\begin{axdef}
  watches : OBS \fun KEY
\where
  watches = \{ theModel \mapsto krows, selObs \mapsto ksel \}
\end{axdef}

A patch step either sets a known field to a valid value or names a value the
setter rejects:

\begin{zed}
  STEP ::= setK \ldata KEY \cross GOODVAL \rdata | badK \ldata KEY \rdata
\end{zed}

\texttt{setK}: set field \texttt{k} to valid value \texttt{g} (the setter
succeeds).  \texttt{badK}: set field \texttt{k} to a rejected value (the setter
raises, triggering rollback).  Each constructor is an injection and their ranges
partition \texttt{STEP}, so the test $head~q \in \ran setK$ distinguishes them.
Two projections recover a step's target and (for a set) its value:

\begin{axdef}
  stepKey : STEP \fun KEY \\
  stepVal : STEP \pfun GOODVAL
\where
  \forall k : KEY; g : GOODVAL @ \\
  \quad~ stepKey~(setK~(k, g)) = k \land stepVal~(setK~(k, g)) = g \\
  \forall k : KEY @ stepKey~(badK~k) = k \\
  \dom stepVal = \ran setK
\end{axdef}

The queue bound keeps every patch (agent-built or re-projection-generated)
finite, and the stack bound caps the reentrancy depth so the carrier is finite:

\begin{axdef}
  maxlen : \nat \\
  maxdepth : \nat
\where
  maxlen = 2 \\
  maxdepth = 3
\end{axdef}

A patch of length two exhibits ``set a watched field, then fail on a second''; a
re-projection writes exactly two keys (\texttt{krows} then \texttt{ksel}).  A
stack of depth two is the deepest the correct design reaches (an outer patch and
its one nested re-projection); depth three gives ProB the room to exhibit the
unbounded self-recursion the guard prevents.

% ============================================================
\section{The Patch Frame Stack}
% ============================================================

Reentrancy means more than one \texttt{apply\_patch} is live at once: the outer
patch is mid-commit, delivering a notification, when the observer starts a nested
patch.  Each live call is a \emph{frame} with its own rollback snapshot, its own
notification buffer, and (once committing) its own set of notifications still to
deliver.  A frame is the Z counterpart of one Python \texttt{apply\_patch} stack
activation.

The stack of frames is modelled as five parallel sequences kept in lockstep, one
per frame attribute, rather than a sequence of a composite frame type --- ProB
enumerates a composite type's unbounded sequence field as an unconstrained
\texttt{ANY}, so the flat encoding is what keeps the search finite.  Frame $i$ is
the tuple $(smode~i,\, squeue~i,\, ssnap~i,\, sprops~i,\, sreproj~i)$, and the
top (innermost live) frame is at index $\# smode$:

\begin{itemize}
  \item \texttt{smode} --- \texttt{setting} while the frame's setters run,
    \texttt{flushing} while its committed notifications are delivered.
  \item \texttt{squeue} --- the frame's unprocessed setters.
  \item \texttt{ssnap} --- its rollback snapshot, a value copy of the fields as
    they were when the frame began.
  \item \texttt{sprops} --- the props this frame has flagged.  While
    \texttt{setting} this is the per-patch notification \emph{buffer} (the model
    of \texttt{self.\_notify\_buffer}, whose reset-on-entry is why each frame's
    buffer is its own); on commit the same set becomes the props still to
    \emph{deliver}, shrinking as each fires.
  \item \texttt{sreproj} --- \texttt{on} for a frame that is the model's own
    re-projection, whose lifecycle raises and lowers the guard.
\end{itemize}

% ============================================================
\section{State}
% ============================================================

The state carries the five stack sequences plus the field values, each
observer's folded value, the last externally-committed dataset, and the guard.
Its predicate carries only the lockstep and length bounds, which are structural
and true in every design; the atomicity property (Section~\ref{sec:inv}) is
deliberately \emph{absent} here, so a violating state stays reachable and ProB
can find it.

\begin{schema}{Sys}
  smode : \seq FMODE \\
  squeue : \seq (\seq STEP) \\
  ssnap : \seq (KEY \fun GOODVAL) \\
  sprops : \seq (\power KEY) \\
  sreproj : \seq SWITCH \\
  build : \seq STEP \\
  fields : KEY \fun GOODVAL \\
  observed : OBS \fun GOODVAL \\
  extern : GOODVAL \\
  guard : SWITCH
\where
  \# smode = \# squeue \land \# smode = \# ssnap \\
  \land \# smode = \# sprops \land \# smode = \# sreproj \\
  \# smode \leq maxdepth \\
  \# build \leq maxlen \\
  \forall i : \dom squeue @ \# (squeue~i) \leq maxlen
\end{schema}

\texttt{build} is the outer patch an agent assembles one step at a time while
the element is idle, before applying it.  Building step-by-step --- rather than
supplying a whole sequence at \texttt{BeginOuter} --- keeps every input a single
\texttt{STEP} over the finite \texttt{STEP} carrier, so ProB enumerates the patch
shapes without trying to enumerate arbitrary sequences.  The element is
\emph{idle} exactly when the stack is empty ($\# smode = 0$).
\texttt{fields} is the element's current field values.  \texttt{observed}
records, for each observer, the value it has folded in --- for \texttt{selObs}
the abstraction of the full selection, for \texttt{theModel} the abstraction of
the model's \emph{dataset}.  \texttt{extern} tracks the last \texttt{krows} value
written by an \emph{external} (non-re-projection) patch: the dataset the agent
actually pushed, which is what the model's dataset must equal.  The field value
and the dataset diverge under a filter --- the element shows the visible subset
while the model holds the full dataset --- which is exactly why the dataset
invariant compares \texttt{observed~theModel} to \texttt{extern}, not to
\texttt{fields~krows}.  \texttt{guard} is the \texttt{\_reprojecting} boolean.

% ============================================================
\section{Initialization}
% ============================================================

\begin{schema}{Init}
  Sys'
\where
  smode' = \langle \rangle \land squeue' = \langle \rangle \\
  \land ssnap' = \langle \rangle \land sprops' = \langle \rangle \\
  \land sreproj' = \langle \rangle \land build' = \langle \rangle \\
  fields' = \{ krows \mapsto a, ksel \mapsto a \} \\
  observed' = \{ theModel \mapsto a, selObs \mapsto a \} \\
  extern' = a \\
  guard' = off
\end{schema}

A single initial state: no patch in progress, every field holds \texttt{a},
every observer has folded \texttt{a}, the last external dataset was \texttt{a},
and the guard is down.

% ============================================================
\section{Beginning a Patch}
% ============================================================

An agent assembles the outer patch one step at a time while idle, appending each
step to \texttt{build} up to the length bound.  Any step is admissible,
including a \texttt{badK} that will later trigger a rollback:

\begin{schema}{Enqueue}
  \Delta Sys \\
  p? : STEP
\where
  \# smode = 0 \\
  \# build < maxlen \\
  build' = build \cat \langle p? \rangle \\
  smode' = smode \land squeue' = squeue \\
  \land ssnap' = ssnap \land sprops' = sprops \land sreproj' = sreproj \\
  fields' = fields \land observed' = observed \\
  \land extern' = extern \land guard' = guard
\end{schema}

\texttt{BeginOuter} applies the assembled patch: it pushes one frame onto each
stack sequence with \texttt{squeue} taking the built steps, takes the rollback
snapshot --- a value copy of the current fields --- and clears \texttt{build}.
The empty patch is admissible; it simply commits with nothing to do:

\begin{schema}{BeginOuter}
  \Delta Sys
\where
  \# smode = 0 \\
  smode' = smode \cat \langle setting \rangle \\
  squeue' = squeue \cat \langle build \rangle \\
  ssnap' = ssnap \cat \langle fields \rangle \\
  sprops' = sprops \cat \langle \emptyset \rangle \\
  sreproj' = sreproj \cat \langle off \rangle \\
  build' = \langle \rangle \\
  fields' = fields \land observed' = observed \\
  \land extern' = extern \land guard' = guard
\end{schema}

% ============================================================
\section{Applying a Patch}
\label{sec:apply}
% ============================================================

Each step processes the head of the top frame's queue.  A successful set
installs the value into the element and \emph{buffers} the notification for that
prop into the top frame's own \texttt{sprops} entry --- it does not fire it:

\begin{schema}{ApplyStepTop}
  \Delta Sys
\where
  \# smode \neq 0 \\
  smode~(\# smode) = setting \land squeue~(\# smode) \neq \langle \rangle \\
  head~(squeue~(\# smode)) \in \ran setK \\
  fields' = fields \oplus \\
  \quad~ \{ stepKey~(head~(squeue~(\# smode))) \mapsto
           stepVal~(head~(squeue~(\# smode))) \} \\
  squeue' = squeue \oplus \{ \# smode \mapsto tail~(squeue~(\# smode)) \} \\
  sprops' = sprops \oplus \\
  \quad~ \{ \# smode \mapsto (sprops~(\# smode) \cup
           \{ stepKey~(head~(squeue~(\# smode))) \}) \} \\
  smode' = smode \land ssnap' = ssnap \land sreproj' = sreproj \\
  observed' = observed \land extern' = extern \\
  \land guard' = guard \land build' = build
\end{schema}

When the top frame's queue is empty it commits: it flips to \texttt{flushing}
and its buffered \texttt{sprops} become its delivery set (the same set, now read
as props-to-deliver --- the model of \texttt{\_notify\_buffer = None} before the
flush loop, with the local delivery list carrying the props on).  An outer patch
that wrote \texttt{krows} records that value as the new external dataset here; a
re-projection frame does not, because its \texttt{krows} write is the filtered
subset, not an agent push:

\begin{schema}{CommitTopExtern}
  \Delta Sys
\where
  \# smode \neq 0 \\
  smode~(\# smode) = setting \land squeue~(\# smode) = \langle \rangle \\
  sreproj~(\# smode) = off \land krows \in sprops~(\# smode) \\
  extern' = fields~krows \\
  smode' = smode \oplus \{ \# smode \mapsto flushing \} \\
  squeue' = squeue \land ssnap' = ssnap \\
  \land sprops' = sprops \land sreproj' = sreproj \\
  fields' = fields \land observed' = observed \\
  \land guard' = guard \land build' = build
\end{schema}

\begin{schema}{CommitTopPlain}
  \Delta Sys
\where
  \# smode \neq 0 \\
  smode~(\# smode) = setting \land squeue~(\# smode) = \langle \rangle \\
  \lnot (sreproj~(\# smode) = off \land krows \in sprops~(\# smode)) \\
  smode' = smode \oplus \{ \# smode \mapsto flushing \} \\
  squeue' = squeue \land ssnap' = ssnap \\
  \land sprops' = sprops \land sreproj' = sreproj \\
  fields' = fields \land observed' = observed \\
  \land extern' = extern \land guard' = guard \land build' = build
\end{schema}

A rejected step aborts the top frame: its fields are restored from its own
snapshot, its buffer is discarded unfired, and the frame is popped off every
stack sequence.  If the aborted frame was a re-projection, popping it lowers the
guard --- the model's \texttt{try/finally} lowers \texttt{\_reprojecting} whether
the nested patch commits or raises.  Because the outer frame beneath is
untouched, its own delivery set survives and it will resume flushing:

\begin{schema}{AbortReproj}
  \Delta Sys
\where
  \# smode \neq 0 \\
  smode~(\# smode) = setting \land squeue~(\# smode) \neq \langle \rangle \\
  head~(squeue~(\# smode)) \in \ran badK \land sreproj~(\# smode) = on \\
  fields' = ssnap~(\# smode) \\
  smode' = front~smode \land squeue' = front~squeue \\
  \land ssnap' = front~ssnap \land sprops' = front~sprops \\
  \land sreproj' = front~sreproj \\
  guard' = off \land observed' = observed \\
  \land extern' = extern \land build' = build
\end{schema}

\begin{schema}{AbortOuter}
  \Delta Sys
\where
  \# smode \neq 0 \\
  smode~(\# smode) = setting \land squeue~(\# smode) \neq \langle \rangle \\
  head~(squeue~(\# smode)) \in \ran badK \land sreproj~(\# smode) = off \\
  fields' = ssnap~(\# smode) \\
  smode' = front~smode \land squeue' = front~squeue \\
  \land ssnap' = front~ssnap \land sprops' = front~sprops \\
  \land sreproj' = front~sreproj \\
  guard' = guard \land observed' = observed \\
  \land extern' = extern \land build' = build
\end{schema}

% ============================================================
\section{Delivering Notifications (the reentrancy)}
\label{sec:deliver}
% ============================================================

A committing frame delivers its props one at a time.  Delivering a \texttt{ksel}
notification folds the selection observer --- a leaf fold with no reentrancy:

\begin{schema}{DeliverSel}
  \Delta Sys
\where
  \# smode \neq 0 \\
  smode~(\# smode) = flushing \land ksel \in sprops~(\# smode) \\
  observed' = observed \oplus \{ selObs \mapsto fields~ksel \} \\
  sprops' = sprops \oplus \{ \# smode \mapsto (sprops~(\# smode) \setminus \{ ksel \}) \} \\
  smode' = smode \land squeue' = squeue \\
  \land ssnap' = ssnap \land sreproj' = sreproj \\
  fields' = fields \land extern' = extern \\
  \land guard' = guard \land build' = build
\end{schema}

Delivering a \texttt{krows} notification is where the model's observer runs.
When the guard is \emph{up}, the notification is the model's own re-projection
write coming back: the observer returns early and folds nothing --- this is the
self-write guard:

\begin{schema}{DeliverRowsSkip}
  \Delta Sys
\where
  \# smode \neq 0 \\
  smode~(\# smode) = flushing \land krows \in sprops~(\# smode) \\
  guard = on \\
  sprops' = sprops \oplus \{ \# smode \mapsto (sprops~(\# smode) \setminus \{ krows \}) \} \\
  smode' = smode \land squeue' = squeue \\
  \land ssnap' = ssnap \land sreproj' = sreproj \\
  observed' = observed \land fields' = fields \\
  \land extern' = extern \land guard' = guard \land build' = build
\end{schema}

When the guard is \emph{down}, the notification is an external dataset refresh.
The model absorbs it (\texttt{observed~theModel} becomes the just-delivered field
value), raises the guard, and pushes a nested re-projection frame that will write
the filtered subset (\texttt{b}) into \texttt{krows} and a projected selection
into \texttt{ksel}.  Crucially the nested frame carries its \emph{own} empty
\texttt{sprops} entry, and the outer frame's entry has already had \texttt{krows}
removed --- the two buffers never share storage:

\begin{schema}{Reproject}
  \Delta Sys
\where
  \# smode \neq 0 \\
  smode~(\# smode) = flushing \land krows \in sprops~(\# smode) \\
  guard = off \land \# smode < maxdepth \\
  observed' = observed \oplus \{ theModel \mapsto fields~krows \} \\
  smode' = smode \cat \langle setting \rangle \\
  squeue' = squeue \cat \langle \langle setK~(krows, b), setK~(ksel, b) \rangle \rangle \\
  ssnap' = ssnap \cat \langle fields \rangle \\
  sprops' = (sprops \oplus \{ \# smode \mapsto
             (sprops~(\# smode) \setminus \{ krows \}) \}) \cat \langle \emptyset \rangle \\
  sreproj' = sreproj \cat \langle on \rangle \\
  guard' = on \land fields' = fields \\
  \land extern' = extern \land build' = build
\end{schema}

The same trigger may instead push a re-projection that \emph{aborts} on its
second key --- the model asking ``can the nested patch fail mid-flight?''  Its
rollback restores the fields to the post-outer-commit values, and the outer frame
resumes flushing its remaining props, so the outer commit is never left
half-delivered:

\begin{schema}{ReprojectAbort}
  \Delta Sys
\where
  \# smode \neq 0 \\
  smode~(\# smode) = flushing \land krows \in sprops~(\# smode) \\
  guard = off \land \# smode < maxdepth \\
  observed' = observed \oplus \{ theModel \mapsto fields~krows \} \\
  smode' = smode \cat \langle setting \rangle \\
  squeue' = squeue \cat \langle \langle setK~(krows, b), badK~ksel \rangle \rangle \\
  ssnap' = ssnap \cat \langle fields \rangle \\
  sprops' = (sprops \oplus \{ \# smode \mapsto
             (sprops~(\# smode) \setminus \{ krows \}) \}) \cat \langle \emptyset \rangle \\
  sreproj' = sreproj \cat \langle on \rangle \\
  guard' = on \land fields' = fields \\
  \land extern' = extern \land build' = build
\end{schema}

When a frame has delivered every prop it is popped off every stack sequence.  A
re-projection frame lowers the guard as it goes; an outer frame leaves the guard
untouched and, being the last frame, returns the element to idle:

\begin{schema}{PopReproj}
  \Delta Sys
\where
  \# smode \neq 0 \\
  smode~(\# smode) = flushing \land sprops~(\# smode) = \emptyset \\
  sreproj~(\# smode) = on \\
  smode' = front~smode \land squeue' = front~squeue \\
  \land ssnap' = front~ssnap \land sprops' = front~sprops \\
  \land sreproj' = front~sreproj \\
  guard' = off \\
  fields' = fields \land observed' = observed \\
  \land extern' = extern \land build' = build
\end{schema}

\begin{schema}{PopOuter}
  \Delta Sys
\where
  \# smode \neq 0 \\
  smode~(\# smode) = flushing \land sprops~(\# smode) = \emptyset \\
  sreproj~(\# smode) = off \\
  smode' = front~smode \land squeue' = front~squeue \\
  \land ssnap' = front~ssnap \land sprops' = front~sprops \\
  \land sreproj' = front~sreproj \\
  guard' = guard \\
  fields' = fields \land observed' = observed \\
  \land extern' = extern \land build' = build
\end{schema}

Every reachable non-idle state enables at least one operation: a \texttt{setting}
top with a non-empty queue enables \texttt{ApplyStepTop} or one of the aborts; a
\texttt{setting} top with an empty queue enables one of the commits; a
\texttt{flushing} top with props to deliver enables \texttt{DeliverSel},
\texttt{DeliverRowsSkip}, or the re-projections; a \texttt{flushing} top with
nothing left enables a pop.  With \texttt{BeginOuter} enabled in every idle
state, the machine never deadlocks in the correct design
(Section~\ref{sec:verify}).

% ============================================================
\section{The Invariant}
\label{sec:inv}
% ============================================================

Two safety clauses must hold whenever the element is idle
($\# smode = 0$).  They are stated as a predicate over \texttt{Sys};
they are \emph{not} placed in the \texttt{Sys} schema, because a predicate there
would be enforced as a guard on every successor, silently disabling any operation
that would break it and masking the very defect this model exists to catch.

\paragraph{Invariant AT --- observable atomicity.}
When no patch is in progress:
\[
  \# smode = 0 \implies
  (observed~selObs = fields~ksel \; \land \; observed~theModel = extern).
\]
The two conjuncts capture the two halves of the contract:

\begin{itemize}
  \item \emph{Selection (\texttt{selObs}).}  After a successful patch the
    selection observer holds the committed \texttt{ksel} value; after a rejected
    patch it holds the pre-patch value alongside a restored field.  Either way
    $observed~selObs = fields~ksel$.  A round-1 or round-2 defect breaks this
    equality --- the observer runs ahead of a rolled-back field, or a mutated
    snapshot fails to restore the field beneath an un-fired observer.
  \item \emph{Dataset (\texttt{theModel}).}  The model's dataset must equal the
    dataset the agent last pushed --- the last \emph{external} \texttt{krows}
    write, \texttt{extern} --- and \emph{not} the filtered subset the model's own
    re-projection wrote back into the field.  With the guard, the model absorbs
    only external writes, so $observed~theModel = extern$ at idle even though
    $fields~krows$ now holds the subset.  Removing the guard breaks this: the
    model folds its own subset in as the dataset.
\end{itemize}

% ============================================================
\section{Verification}
\label{sec:verify}
% ============================================================

Type-checking is a \texttt{fuzz} pass:
\begin{verbatim}
  fuzz -t docs/patch_atomicity.tex
\end{verbatim}

Invariant AT is checked by asking ProB whether its negation is \emph{reachable}.
A goal that is not found means the property holds on every reachable state:
\begin{verbatim}
  # Invariant AT (must report: goal NOT found)
  probcli docs/patch_atomicity.tex -model_check -nodead \
    -goal "size(smode)=0 & (observed(selObs)/=fields(ksel)
           or observed(theModel)/=extern)"
\end{verbatim}

Deadlock-freedom is the state-space search with no goal; the machine is kept live
by \texttt{BeginOuter} in every idle state and by an enabled operation in every
in-progress state, so a correct design never deadlocks:
\begin{verbatim}
  # Deadlock-freedom (must report: no deadlock, no error)
  probcli docs/patch_atomicity.tex -model_check
\end{verbatim}

The carriers are concrete free types, so both searches are already exhaustive
--- there is no \texttt{DEFAULT\_SETSIZE} to raise.  Robustness against the
bounds is checked by re-running with \texttt{maxdepth} lowered to \texttt{2}
(the correct design's true maximum depth): the invariant still holds and no
deadlock appears, confirming the depth-three headroom is only there to give the
buggy variant room to recurse.

% ============================================================
\section{Fidelity: reproducing the three defects}
\label{sec:fidelity}
% ============================================================

A model that cannot reproduce the known defects proves nothing.  Each variant is
a single-schema edit; the model checker is re-run, and the named goal must now be
\emph{found} (or, for the guard, a deadlock must appear).

\paragraph{Round 2 --- observer notified inside the patch.}
Model eager notification: in \texttt{ApplyStepTop}, instead of adding the prop to
the frame's \texttt{sprops} buffer, fold it immediately into its observers ---
$observed' = observed \oplus \{ o : OBS | watches~o = k @ o \mapsto g \}$ for the
just-installed key $k$ and value $g$, leaving \texttt{sprops} unchanged so
nothing is buffered.  Now a patch that sets \texttt{ksel} and then hits a
\texttt{badK} step folds the intermediate selection into \texttt{selObs},
\texttt{AbortOuter} restores the field but cannot un-fold the observer, and at
the resulting idle state $observed~selObs \neq fields~ksel$: the AT goal is
\emph{found}.

\paragraph{Round 1 --- snapshot aliases a mutable model.}
Model the shallow snapshot capturing a shared reference: in \texttt{ApplyStepTop}
mutate the top frame's snapshot alongside the field ---
$ssnap' = ssnap \oplus \{ \# smode \mapsto (ssnap~(\# smode) \oplus \{ k \mapsto g \}) \}$
--- leaving the (correctly deferred) notification alone.  Now a patch that sets
\texttt{ksel} then hits a \texttt{badK} reaches \texttt{AbortOuter}, which
restores $fields' = ssnap~(\# smode)$ --- but the snapshot was mutated with the
field, so the field is \emph{not} restored, while the deferred \texttt{selObs}
still holds the pre-patch value.  At idle $fields~ksel \neq observed~selObs$: the
AT goal is \emph{found} from the opposite side.

\paragraph{Round 3 --- the self-write guard removed.}
Model the missing \texttt{\_reprojecting} guard: delete the \texttt{DeliverRowsSkip}
schema and drop the \texttt{guard = off} conjunct from \texttt{Reproject} and
\texttt{ReprojectAbort}, so \emph{every} \texttt{krows} delivery re-absorbs and
re-projects.  Now the nested re-projection's own \texttt{krows} write, delivered
during \emph{its} commit, is folded straight back in as the dataset ---
$observed~theModel$ becomes the filtered subset \texttt{b} while the agent pushed
\texttt{a} --- and it pushes yet another re-projection, and so on until the depth
bound.  Two signals witness the defect:
\begin{verbatim}
  # the observer folds its own subset in as the dataset
  probcli docs/patch_atomicity.tex -model_check -nodead \
    -goal "guard=on & observed(theModel)/=extern"       (* found *)

  # the self-recursion, bounded only by the state space
  probcli docs/patch_atomicity.tex -model_check          (* deadlock at maxdepth *)
\end{verbatim}
The first goal is unreachable in the correct design --- while the guard is up the
model has already absorbed the external dataset and folds nothing more --- so its
reachability is precisely the guard's absence.  The deadlock is the recursion
hitting \texttt{maxdepth} with a \texttt{krows} still undelivered and no frame it
is allowed to push: the finite image of an unbounded re-entrancy.

\paragraph{Under the correct design none of the three traces exists.}
\texttt{ApplyStepTop} defers every notification into the frame's own buffer and
leaves the snapshot untouched; \texttt{AbortReproj} / \texttt{AbortOuter} restore
a value snapshot and fire nothing; each commit hands its buffer to its own
delivery set; and \texttt{DeliverRowsSkip} stops the model re-absorbing its own
re-projection.  The AT search returns \emph{goal not found}, the guard goal is
unreachable, and the machine is deadlock-free.

% ============================================================
\section{Nested-case partitions and test coverage}
\label{sec:partition}
% ============================================================

The reentrancy adds four partitions to the single-patch atomicity partitions the
earlier rounds already covered.  Each is a reachable class of trace in the model;
the table names the operation sequence, the AT outcome, and the test in
\texttt{tests/test\_table\_composition.py} that exercises it.

\begin{center}
\begin{tabular}{p{4.6cm} p{3.0cm} p{5.2cm}}
\hline
\textbf{Partition (model trace)} & \textbf{AT outcome} & \textbf{Covering test} \\
\hline
External rows refresh under a filter: \texttt{BeginOuter}(rows) $\to$ commit $\to$
\texttt{Reproject} $\to$ nested writes subset $\to$ \texttt{DeliverRowsSkip} (guard
up) $\to$ pop. &
$observed~theModel = extern$ (full dataset, not subset). &
\texttt{test\_agent\_rows\_patch\_refreshes\_the\_dataset\_under\_a\_filter} ---
the agent push becomes the dataset and a later filter uses it. \\
\hline
Self-write guard holds: the nested re-projection's own \texttt{krows} delivery is
skipped, not re-absorbed. &
Dataset stays the full set; guard goal unreachable. &
\texttt{test\_the\_models\_own\_reprojection\_keeps\_the\_dataset} --- a
filter-clear restores every row. \\
\hline
Dataset change reconciles the full selection (ids gone from the data leave it),
distinct from a filter hiding them. &
Selection folds to the reconciled set. &
\texttt{test\_agent\_rows\_patch\_reconciles\_the\_full\_selection}. \\
\hline
Dual rows+selection outer patch under a filter, both props buffered and delivered,
one triggering the nested re-projection. &
Both \texttt{observed} clauses hold. &
\texttt{test\_dual\_rows\_and\_selection\_patch\_keeps\_the\_selection}. \\
\hline
Nested re-projection \emph{aborts} mid-flight (\texttt{ReprojectAbort}); the outer
commit still delivers its remaining props. &
Outer never half-delivered; both clauses hold. &
\emph{Gap} --- covered in the model, not in a test. \\
\hline
\end{tabular}
\end{center}

The single-patch atomicity partitions remain covered by
\texttt{test\_a\_failed\_multi\_key\_patch\_leaves\_the\_model\_unchanged} and
\texttt{test\_a\_successful\_patch\_still\_folds\_through\_to\_the\_model} in
\texttt{TestWriteThroughAtomicity}.

\paragraph{The one gap.}  The \texttt{ReprojectAbort} partition --- a nested
re-projection that raises after writing its first key --- is proven safe by the
model (the outer commit finishes delivering, and both AT clauses hold at idle)
but has no dedicated test.  In the shipped code a re-projection writes only
validated subset rows and a projected selection, so it does not raise in
practice; the partition guards against a future setter that could.  A regression
test that forces the model's inner \texttt{apply\_patch} to reject its second key
and asserts the outer selection notification still lands would close it.  The
model stands as the regression artefact for the interleaving until then.

\end{document}
